\section*{Chapitre \ref{ch:m3:getting-access-crypto-venues} --- Accès aux marchés~: plateformes crypto}

\begin{solution}{exo:m3:getting-access-crypto-venues:1}
USD~100 millions est au palier VIP~4 (0.040~\%/0.052~\%)~: $100\text{ millions} \times (0.7 \times 4 + 0.3 \times 5.2)$ points de base, soit USD~43\,600.
\end{solution}

\begin{solution}{exo:m3:getting-access-crypto-venues:2}
Identité et constitution en société, bénéficiaires effectifs et dirigeants, origine des fonds, agréments et statut réglementaire, stratégie de trading. La
plateforme supporte seule le risque juridique de servir un client sanctionné ou non agréé, ou de blanchir des fonds, et ses propres agréments dépendent de ces contrôles.
\end{solution}

\begin{solution}{exo:m3:getting-access-crypto-venues:3}
Dix par seconde pendant dix secondes font 100, mais 200\,000 par jour ne représentent que 2.3 par seconde en moyenne~: une stratégie qui cote
toute la journée se heurte à la règle quotidienne bien avant de saturer la règle courte.
\end{solution}

\begin{solution}{exo:m3:getting-access-crypto-venues:4}
USD~11\,040 par mois~: de 3.6 points de base au barème propre de la plateforme C à 2.68 au palier Binance équivalent.
\end{solution}

\begin{solution}{exo:m3:getting-access-crypto-venues:5}
USD~10\,500 par mois~; rater le seuil d'un seul point fait perdre la totalité, puisque la ristourne n'est versée que si l'exigence est remplie.
\end{solution}

\begin{solution}{exo:m3:getting-access-crypto-venues:6}
Résilience et information~: une plateforme peut faire défaut, suspendre les retraits ou tomber en panne pendant une crise, et les prix et la liquidité de chaque plateforme éclairent la cotation
sur les autres~; de petites allocations entretiennent aussi les comptes, les connexions et les relations.
\end{solution}

\begin{solution}{exo:m3:getting-access-crypto-venues:7}
Environ 64.0~\% à 10 points de base et 93.6~\% à 20~: à 10, la limite sur la fourchette s'impose bien avant la règle de profondeur.
\end{solution}

\begin{solution}{exo:m3:getting-access-crypto-venues:8}
Négocier entre ses propres \omterm{def:m3:centralised-exchanges:subaccount}{sous-comptes} constitue du \omterm{def:m3:spot-markets:wash}{wash trading}~; les plateformes l'excluent des paliers lorsqu'elles peuvent le détecter, il enfreint leurs conditions et peut constituer une manipulation illégale.
Il paie aussi des frais preneur sur la moitié de ce volume.
\end{solution}

\begin{solution}{pb:m3:getting-access-crypto-venues:1}
\textbf{1.} USD~52\,320 (Binance), 28\,800 (Kraken), 43\,200 (C), 86\,400 (D) et 110\,400 (E)~: USD~321\,120. \textbf{2.} USD~80\,400 et
63\,000~: USD~143\,400. \textbf{3.} USD~177\,720 par mois. \textbf{4.} Binance passe du VIP~4 au VIP~5 et Kraken du palier de USD~100 millions au
palier de USD~250 millions~; la commission de maker de Kraken est nulle à partir de USD~10 millions, de sorte que 70~\% de son volume ne paie rien. \textbf{5.} Cela réduirait le coût de C de
USD~11\,040, ce qui rend l'allocation à trois plateformes moins chère que la répartition et plus proche de la concentration. \textbf{6.} $300\text{ millions} \times 70\% \times
0.5$ points de base~: USD~10\,500 par mois. \textbf{7.} 87.3~\%~: la stratégie ne remplit pas les conditions. \textbf{8.} Maintenir la pleine taille jusqu'à 2.25 fois la volatilité normale
au lieu de 2.0~: 91.0~\%. \textbf{9.} Cotations à pleine taille sur davantage de minutes volatiles, lorsque la sélection adverse est maximale~; la ristourne doit dépasser les
pertes attendues supplémentaires. \textbf{10.} Maintenir des cotations conformes sur les marchés rapides exige des mises à jour, que les règles sur le nombre d'ordres rationnent~; la société doit budgétiser
ses messages pour les moments que mesure le programme. \textbf{11.} La concentration fait passer USD~300 millions par mois sur chacune des deux plateformes, et
le collatéral et les soldes qui s'y trouvent~; la répartition réduit de moitié chaque exposition, mais en ajoute trois. \textbf{12.} Le collatéral détenu auprès d'un dépositaire et répliqué sur les plateformes
réduit ce que peut emporter chaque défaillance de plateforme, faisant de la concentration des négociations une moindre concentration de crédit. \textbf{13.} La liquidité et les prix sur
ces plateformes, un repli si l'une des deux fait défaut, et l'option d'y négocier plus tard. \textbf{14.} Par le coût réalisé de ses propres transactions sur chaque plateforme
(fourchette payée, slippage, markouts), et non par les seules commissions. \textbf{15.} Mensuellement, et à chaque changement de barème, de programme ou de volume. \textbf{16.}
Cela réduit les coûts de la société et approfondit deux carnets d'ordres, mais il amincit les autres et concentre la liquidité et le risque du marché. \textbf{17.} Les ristournes achètent
le comportement de cotation que la plateforme attend des sociétés capables de le fournir~; réduire les frais pour tous n'achète rien de précis. \textbf{18.} Peu~: deux ou
trois là où elles peuvent atteindre de bons paliers, plus une plateforme maintenue petite pour la résilience. \textbf{19.} \emph{Résultat à retenir~:} USD~177\,720 par mois économisés en se concentrant
sur deux plateformes~; la ristourne suppose une \omterm{def:m3:getting-access-crypto-venues:uptime}{exigence de disponibilité} de 90~\%, que la stratégie n'atteint (91.0~\%) que si elle conserve la pleine taille jusqu'à 2.25 fois la volatilité normale,
contre 87.3~\% en l'état. \textbf{20.} Accès aux flux d'ordres de ses clients, à des conditions qui s'améliorent avec le volume et le risque que le teneur de marché apporte.
\end{solution}

\begin{solution}{iq:m3:getting-access-crypto-venues:1}
Les frais diminuent avec le volume sur chaque plateforme, de sorte que concentrer le volume abaisse les frais moyens~; l'allocation la moins chère place l'essentiel du volume là où le barème le récompense
sous réserve de la résilience, de la liquidité et des limites de crédit par plateforme.
\iqlookfor{effets marginaux des paliers et contraintes hors frais.}
\end{solution}

\begin{solution}{iq:m3:getting-access-crypto-venues:2}
Ses agréments et sa juridiction, son actionnariat, ses états financiers ou sa \omterm{def:m3:centralised-exchanges:por}{preuve de réserves} et de passifs, les modalités de conservation et de séparation des actifs, le bilan des
retraits, son assurance, le traitement des actifs clients en cas d'insolvabilité, ainsi que la limite de crédit de la société et sa politique de balayage à son égard.
\iqlookfor{analyse de crédit de la plateforme.}
\end{solution}

\begin{solution}{iq:m3:getting-access-crypto-venues:3}
Mesurer la disponibilité en continu, comme le fait la plateforme~; donner la priorité aux mises à jour de cotations sur les instruments exposés au risque de défaillance~; regrouper les mises à jour~; garder un budget de messages
pour les marchés rapides~; élargir dans la fourchette autorisée plutôt qu'annuler~; alerter avant que la mesure mensuelle ne descende sous le seuil.
\iqlookfor{mesure, priorisation et budget de messages.}
\end{solution}

\begin{solution}{iq:m3:getting-access-crypto-venues:4}
Seulement si la ristourne dépasse les pertes attendues de sélection adverse liées à la cotation pendant ces périodes, mesurées à partir des markouts propres à la société dans les
minutes volatiles~; la réponse varie selon l'instrument.
\iqlookfor{ristourne comparée aux pertes extrêmes mesurées.}
\end{solution}

\begin{solution}{iq:m3:getting-access-crypto-venues:5}
Comparer son volume à sa \omterm{def:m3:spot-markets:depth}{profondeur affichée} et à son activité sur le web ou on-chain, tester la taille de ses transactions au regard des arrondis et de la loi de Benford
(voir le \cref{ch:m3:spot-markets}), vérifier la co-variation de ses prix avec celle des principales plateformes, et négocier de petits montants pour voir la profondeur réelle.
\iqlookfor{preuve indépendante d'une activité réelle.}
\end{solution}

\begin{solution}{iq:m3:getting-access-crypto-venues:6}
Pour chaque plateforme~: volume et palier, frais et ristournes perçus, disponibilité au regard de chaque programme, usage des messages au regard des limites, soldes et collatéral
au regard des limites de crédit, incidents de retrait et de connectivité, et qualité d'exécution~; avec des alertes en cas de franchissement de seuil et une recommandation mensuelle sur
l'allocation.
\iqlookfor{coûts, obligations, limites et crédit dans une seule vue.}
\end{solution}
